% ===========================================================================
%  附录 D：MPC
% ===========================================================================

\section{模型预测控制的约束接口}
\label{sec:mpc}

模型预测控制（MPC）能够在有限预测时域内显式处理摆角、差速、转速和状态约束，
适合研究本构型的饱和管理和过渡轨迹；其效果依赖预测模型、状态估计和实时求解器
\cite{mayne2000constrained,rawlings2020model,kang2009control}。

\subsection{离散预测模型}

以状态 $\bm x=[\bm p_N^T,\bm v_P^T,\bm q^T,\bm\omega^T,
\omega_f,\omega_t]^T$ 和输入
$\bm u=[\omega_0,\Delta\omega,\delta_t,\delta_f]^T$ 为例，
离散模型应由\cref{eq:newton_body,eq:euler_body,eq:quaternion_kinematics,eq:motor_lag}
共同构成：
\begin{equation}
\bm x_{k+1}=\bm f_d(\bm x_k,\bm u_k;\bm p).
\label{eq:mpc_dynamics}
\end{equation}
式\eqref{eq:mpc_dynamics}必须保留
$\bm\omega\times(\bm I\bm\omega)$、$\bm\omega\times\bm h_r$、
电机一阶动态和四元数单位范数处理。只用 $\bm M=\bm I\bm\nu$ 再次除以惯量，
或忽略 $\tau_m$ 计算控制带宽，都会产生量纲或时间尺度错误。

\subsection{有限时域优化}

典型目标可写为
\begin{equation}
\min_{\bm u_{0:N-1}}
\sum_{k=0}^{N-1}
\left(\|\bm e_{R,k}\|_{\bm Q_R}^2
+\|\bm\omega_k-\bm\omega_{d,k}\|_{\bm Q_\omega}^2
+\|\Delta\bm u_k\|_{\bm R}^2\right)
+\|\bm x_N-\bm x_{d,N}\|_{\bm P}^2,
\label{eq:mpc_ocp}
\end{equation}
约束至少包括固件 $\pm15^\circ$ 摆角、$\pm0.7$ 差速、
$0\leq\omega_i\leq\omega_{\max}$、输入变化率、最低推力和四元数范数。
若研究转捩，还需加入高度、空速、迎角和推力裕度约束。

线性时变 MPC 可在每拍围绕当前工作点离散化；非线性 MPC 则直接使用
式\eqref{eq:mpc_dynamics}。\texttt{acados} 提供面向嵌入式最优控制的代码生成框架
\cite{verschueren2022acados}，\texttt{TinyMPC} 展示了资源受限微控制器上的凸 MPC 实现
\cite{cassel2024tinympc}。这些工具说明实时求解具有可行路径，但不能替代本项目的
最坏执行时间、数值失败回退和闭环稳定性验证。

\subsection{当前结论边界}

当前工程没有与固件逐式对齐、带运行清单和最坏时延测试的 MPC 实现，
因此本文不报告 MPC 性能数字，也不称其为已完成控制器。其最现实的近期用途是作为离线
约束轨迹和分配基准，与级联控制器在相同模型和限制下比较；只有在求解失败策略、
HIL 时序和实机安全门控完成后，才适合进入在线飞行验证。
